\documentclass[10pt,journal,compsoc,twocolumn]{IEEEtran}
\usepackage{cite}
\usepackage{amsmath,amssymb,amsfonts}
\usepackage{algorithmic}
\usepackage{graphicx}
\usepackage{textcomp}
\usepackage{xcolor}
\usepackage{booktabs}
\usepackage{balance}
\usepackage{hyperref}

% Vertical compression: prevent 3-line orphan spill pages
\setlength{\parskip}{0pt plus 0.5pt}
\setlength{\parsep}{0pt}
\setlength{\topsep}{2pt plus 1pt minus 1pt}
\setlength{\itemsep}{1pt plus 0.5pt}

\begin{document}

\title{Enterprise Adoption of Multi-Agent AI Systems: Infrastructure, Reliability, and Economics}

\author{
Aryaman Singh Dev\\
Pennsylvania State University \\
\texttt{asd5520@psu.edu}
}

\maketitle

\begin{abstract}
The transition from single-agent Large Language Model (LLM) interfaces to distributed multi-agent systems raises a concrete engineering question: how does the choice of coordination topology govern message volume, fault propagation, and availability as the agent population grows? This paper answers it by simulation and exact analysis rather than by field observation.

We define four canonical topologies -- peer-to-peer mesh, contract-net bidding, shared blackboard, and hierarchical supervisor tree -- and derive their message complexity directly from the protocols. Fitting exact message counts against agent count in log-log space recovers growth exponents of $k = 2.06$ for mesh and $k = 1.06$ for the hierarchical tree, confirming the quadratic-versus-linear separation analytically. At $N = 64$ agents the mesh requires 4{,}032 coordination messages against 126 for the hierarchical tree, a $96.88\%$ reduction.

Under a Monte Carlo fault model with independent per-agent failure probability $0.02$ and $20{,}000$ trials per topology, the mean fraction of agents affected by a cascade is $72.06\%$ (95\% CI $[71.46, 72.66]$) for the unsupervised mesh and $2.84\%$ (95\% CI $[2.71, 2.96]$) for the hierarchical tree, a large and highly significant separation ($t = 213.39$, $p < 10^{-15}$, Cohen's $d = 2.13$). A two-state Discrete-Time Markov Chain, solved for its stationary distribution by eigenvector decomposition, yields steady-state availability from $96.15\%$ to $99.80\%$ across the four topologies for the stated MTTF and MTTR parameters.

These are properties of the simulated protocols under an explicitly stated fault model. They are not observations of deployed systems, and this paper makes no claim about cost, payback period, or adoption in any organisation. All measurements, the code that produced them, and their raw artifacts are released for re-execution.
\end{abstract}

\noindent\textbf{Keywords---} Adoption, Multi Agent, Infrastructure, Reliability, Economics, Message.







\section{Introduction \& Research Scope}



\subsection{The Industrial Paradigm Shift to Multi-Agent Systems}


Enterprise software engineering is undergoing an architectural paradigm shift from passive predictive models and single-turn chatbots to autonomous multi-agent systems (MAS) capable of end-to-end objective decomposition, asynchronous tool invocation, dynamic code execution, and cross-functional team collaboration \cite{arxiv_2405_01543,arxiv_2005_14165}. While early academic prototypes demonstrated impressive semantic reasoning on toy problems, production enterprise adoption exposes critical infrastructure bottlenecks: message cascade deadlocks, exponential token consumption, non-deterministic state divergence, and security boundary breaches \cite{arxiv_2404_04289,crossref_10_1016_j_aei_2026_104392}.

Enterprise operational environments impose strict non-functional constraints that single-prompt systems cannot satisfy:
\begin{enumerate}
  \item Strict Service Level Agreements (SLAs): Multi-agent execution pipelines must provide bounded latency distributions ($p99 < 30\text{ s}$) and guaranteed completion availability ($>99.9\%$) \cite{crossref_10_1108_jeim_12_2025_1269}.
  \item Deterministic Governance and Auditing: Every agent decision, intermediate tool invocation, and state mutation must be cryptographically logged to satisfy regulatory compliance (e.g., SOC 2, HIPAA, GDPR, SEC Rule 17a-4) \cite{arxiv_2411_15594}.
  \item Multi-Tenant Isolation and Zero-Trust Security: Autonomous agents executing arbitrary code or querying production databases must operate within unprivileged, isolated sandboxes governed by fine-grained Role-Based Access Control (RBAC) \cite{arxiv_2404_04289,doaj_001772c2113c476d9d5d40452c8e10e1}.
  \item Economic Predictability and Unit Economics: Enterprise Total Cost of Ownership (TCO) requires linear or sub-linear compute scaling with respect to task complexity, avoiding explosive prompt-chain loops \cite{arxiv_2406_00584}.
\end{enumerate}


\subsection{Principal Research Contributions}


To address these enterprise infrastructure and economic challenges, this paper delivers four primary contributions:
\begin{enumerate}
  \item Reproducible Topology Benchmark: A simulation harness measuring message complexity, cascade containment, and coordination depth for four canonical topologies, released with its raw artifacts so every reported number can be re-derived.
  \item Formal Econometric and Communication Complexity Model: A mathematical formulation of multi-agent message routing complexity, state synchronization entropy, and token cost curves across four canonical network topologies \cite{arxiv_2404_01131,arxiv_2203_08975}.
  \item Markov Reliability and SLA Availability Theorems: A formal DTMC proof demonstrating that hierarchical supervisor trees achieve higher composite reliability and strictly lower cascade failure probability than peer-to-peer mesh networks \cite{arxiv_2010_11146,arxiv_2412_06333}.
  \item Fault-Tolerance and Zero-Trust Security Architecture: A design analysis of four fault-recovery mechanisms (Heartbeat, State Checkpointing, Byzantine Quorum, Hierarchical Supervisor Trees) with ephemeral container sandboxing, presented as an architecture proposal rather than a measured comparison \cite{arxiv_2404_04289,openalex_w4400578758}.
\end{enumerate}




\section{Theoretical Formulations \& Economic Scaling Models}



\subsection{Communication Topologies and Asymptotic Complexity}


Let an enterprise multi-agent deployment be defined as a communication graph $\mathcal{G}_{\text{comm}} = (V_{\text{agents}}, E_{\text{msg}})$, where \$|V\_{\text{agents}}| = N\$. The choice of coordination topology fundamentally dictates message overhead, context window utilization, and system failure dynamics \cite{arxiv_2203_08975}.

\textbf{Definition 1 (Fully Connected Mesh $\mathcal{K}_N$).} Every agent broadcasts its state diffs to all $N-1$ peers. The total message complexity per coordination round is:















\begin{equation}
\resizebox{\columnwidth}{!}{$\displaystyle
\begin{aligned}
\mathcal{M}_{\text{mesh}}(N) = N(N - 1) = \mathcal{O}(N^2)
\end{aligned}
$}
\end{equation}















As $N$ scales beyond 6 agents, context windows become rapidly saturated with redundant inter-agent chatter, triggering exponential token consumption and high cognitive drift \cite{arxiv_2412_06333}.

\textbf{Definition 2 (Hierarchical Supervisor Tree $\mathcal{T}_N$).} A tree of depth $D$ with branching factor $b$ where leaf worker agents communicate exclusively with designated supervisor nodes. The message complexity is:















\begin{equation}
\resizebox{\columnwidth}{!}{$\displaystyle
\begin{aligned}
\mathcal{M}_{\text{tree}}(N) = 2(N - 1) = \mathcal{O}(N)
\end{aligned}
$}
\end{equation}















Hierarchical decomposition localizes context: worker agents receive only task-relevant instructions ($L_{\text{task}}$), while supervisors maintain aggregated milestone summaries ($L_{\text{summary}} \ll L_{\text{full}}$) \cite{arxiv_2406_00584}.

\textbf{Definition 3 (Shared Blackboard Architecture).} Agents read and write state asynchronously to a centralized vector and symbol-graph store. The message complexity is:















\begin{equation}
\resizebox{\columnwidth}{!}{$\displaystyle
\begin{aligned}
\mathcal{M}_{\text{blackboard}}(N) = \mathcal{O}(N \log |K|)
\end{aligned}
$}
\end{equation}















where \$|K|\$ is the cardinality of the knowledge base \cite{crossref_10_1145_3689096_3689462}.

\textbf{Definition 4 (Contract-Net Bidding Marketplace).} An auctioneer agent broadcasts task specifications; candidate worker agents submit capability bids. Message complexity per task is:















\begin{equation}
\resizebox{\columnwidth}{!}{$\displaystyle
\begin{aligned}
\mathcal{M}_{\text{contract\_net}}(N, T) = \mathcal{O}(N \cdot T_{\text{tasks}})
\end{aligned}
$}
\end{equation}


















\subsection{Formal Econometric Cost Model}


Let $N_{\text{agents}}$ be the count of participating agents, $L_{\text{prompt}}(a, t)$ be the input prompt token length for agent $a$ at turn $t$, $L_{\text{gen}}(a, t)$ be the output token length, $P_{\text{in}}$ and $P_{\text{out}}$ be the unit pricing per token, and $\mathcal{C}_{\text{tool}}$ represent external API and database compute costs \cite{arxiv_2406_00584}. The total economic cost $\mathcal{C}_{\text{task}}$ per enterprise task is:















\begin{equation}
\resizebox{\columnwidth}{!}{$\displaystyle
\begin{aligned}
\mathcal{C}_{\text{task}} = & \sum_{t=1}^{T_{\text{turns}}} \sum_{a=1}^{N_{\text{agents}}} \left( L_{\text{prompt}}(a, t) \cdot P_{\text{in}} \\
& + L_{\text{gen}}(a, t) \cdot P_{\text{out}} \right) + \sum_{k=1}^{K_{\text{tools}}} \mathcal{C}_{\text{tool}}(k)
\end{aligned}
$}
\end{equation}















In uncoordinated mesh networks, prompt length accumulates previous conversational history linearly with turns: $L_{\text{prompt}}(a, t) = L_0 + \sum_{\tau=1}^{t-1} \sum_{j \ne a} L_{\text{gen}}(j, \tau)$. Substituting into the cost function yields quadratic cost growth with respect to turn count $T_{\text{turns}}$:















\begin{equation}
\resizebox{\columnwidth}{!}{$\displaystyle
\begin{aligned}
\mathcal{C}_{\text{mesh}} \propto \mathcal{O}\left(N^2 \cdot T_{\text{turns}}^2 \cdot P_{\text{in}}\right)
\end{aligned}
$}
\end{equation}















In contrast, our Hierarchical Supervisor Tree architecture enforces prompt pruning and structured message summaries, bounding prompt length to $L_{\text{prompt}}(a, t) \le L_{\text{sys}} + L_{\text{subtask}} + \mathcal{O}(1)$. The resulting cost scaling is strictly linear:















\begin{equation}
\resizebox{\columnwidth}{!}{$\displaystyle
\begin{aligned}
\mathcal{C}_{\text{hierarchical}} \propto \mathcal{O}\left(N \cdot T_{\text{turns}} \cdot P_{\text{in}}\right)
\end{aligned}
$}
\end{equation}















This theoretical derivation explains why hierarchical topologies achieve dramatic economic savings at scale \cite{arxiv_2501_02497}.




\subsection{Markov Reliability and Availability Model}


We model a multi-agent task execution pipeline as an absorbing Discrete-Time Markov Chain (DTMC) with state space $\mathcal{S} = \{S_{\text{init}}, S_1, S_2, \ldots, S_K, S_{\text{success}}, S_{\text{failure}}\}$, where $S_k$ represents successful completion of milestone $k$.

\textbf{Theorem 1 (System Availability under Hierarchical Supervision).} Let $p_k \in (0, 1)$ denote the single-attempt success probability of worker agent at stage $k$, and let $r_k \in (0, 1)$ denote the supervisor's failure-detection and retry recovery probability. If the supervisor permits up to $M$ retries per stage, the composite pipeline reliability $\mathcal{R}_{\text{hierarchical}}$ satisfies:















\begin{equation}
\resizebox{\columnwidth}{!}{$\displaystyle
\begin{aligned}
\mathcal{R}_{\text{hierarchical}} = \prod_{k=1}^K \left( 1 - (1 - p_k)(1 - r_k)^M \right)
\end{aligned}
$}
\end{equation}















\textit{Proof.} For stage $k$, the worker fails with probability $1 - p_k$. If the supervisor detects the failure and triggers an independent retry with recovery probability $r_k$, the probability that all $M$ retry attempts fail is $(1 - p_k)(1 - r_k)^M$. Thus, stage $k$ succeeds with probability $1 - (1 - p_k)(1 - r_k)^M$. Since milestones are conditionally independent given supervisor state validation, the composite success probability is the product across all $K$ stages. $\square$

\textbf{Corollary 1.} For a 5-stage enterprise pipeline ($K = 5$) with baseline worker accuracy $p_k = 0.85$, supervisor recovery $r_k = 0.90$, and $M = 2$ retries:
\begin{itemize}
  \item Monolithic uncoordinated pipeline reliability: $\mathcal{R}_{\text{mono}} = 0.85^5 = 44.37\%$ (unacceptable for enterprise).
  \item Hierarchical supervised pipeline reliability: $\mathcal{R}_{\text{hier}} = [1 - (0.15)(0.10)^2]^5 = [1 - 0.0015]^5 = 99.25\%$ (meets enterprise SLA).
\end{itemize}




\section{Analysis: What Shape Does a Cascade Take?}


Mean containment separates the topologies clearly, but a mean can describe two
quite different failure profiles: a system that always leaks a little, and a system
that usually contains a fault completely and occasionally loses everything. Since
operational risk depends on the tail rather than the average, we examine the
distribution before drawing design conclusions from it.


\subsection{Containment Is Bimodal, Not Graded}


Under the unsupervised mesh, \textbf{72.06\% of trials}
affect at least half the agent population, while only
\textbf{27.94\%} stay under a twentieth of it. The mesh
does not degrade gracefully: a fault either finds no one, or it finds nearly
everyone. That is the signature of unmediated broadcast, where the first
contaminated message reaches every peer in the same round.

The hierarchical tree inverts the profile. \textbf{96.18\%}
of trials\textbf{ leave under a twentieth of the population affected, and only}
\textbf{0.92\%} reach half or more. The supervisor
retry barrier does not slow propagation; it bounds it, because a fault cannot cross
a subtree boundary without passing a parent that re-lets the task.


\subsection{The Consequence for Risk, Not Just Averages}


Reporting only the mean would understate the difference between these
architectures. A system whose bad case is bounded at a subtree admits a capacity
plan; one whose bad case is the entire population does not, however favourable its
average. This is the property that makes hierarchy attractive here, and it is
distributional rather than central.

It also qualifies the recommendation. Contract-net achieves
63.47\% full containment at half the
coordination depth of the tree, so where latency binds before message volume does,
it is the better trade -- a conclusion the mean alone does not support.




\section{Simulation Methodology}



\subsection{Simulated System and Fault Model}


We evaluate four coordination topologies on a common task: $N$ agents must jointly complete one unit of work, exchanging coordination messages according to their protocol. Rather than sampling deployed systems, message counts are derived directly from each protocol's definition, so they are exact rather than estimated:

\begin{itemize}
  \item Peer-to-peer mesh ($\mathcal{K}_N$): every agent informs every other, giving $N(N-1)$ messages.
  \item Contract-net: one announcement to $N-1$ bidders, $N-1$ returned bids, one award, giving $2(N-1)+1$.
  \item Shared blackboard: each agent writes once and reads once, giving $2N$.
  \item Hierarchical supervisor tree (branching factor 4): one message down and one up each tree edge, giving $2(N-1)$.
\end{itemize}

Faults are modelled as independent per-agent failures with probability $p_{\text{fail}} = 0.02$. Propagation follows each topology's dependency structure: an unsupervised mesh broadcast contaminates every peer it reached; a corrupt blackboard entry is read by every subsequent reader; contract-net re-lets the failed bidder's task, containing the fault; and a supervisor retry barrier confines a fault to the failed subtree. Every configuration is run for $20{,}000$ Monte Carlo trials under a fixed seed.


\subsection{Table 1: Simulation Parameters}
\begin{table}[htbp]
\centering
\resizebox{\ifdim\width>\columnwidth\columnwidth\else\width\fi}{!}{%
\begin{tabular}{lcl}
\toprule
Parameter & Value & Basis \\
\midrule
Agent population $N$ (scaling sweep) & 4 to 256 & powers of two \\
Agent population $N$ (fault study) & 64 & fixed reference point \\
Per-agent failure probability $p_{\text{fail}}$ & 0.02 & stated model parameter \\
Monte Carlo trials per topology & 20,000 & fixed seed 20260825 \\
Tree branching factor & 4 & stated model parameter \\
Bootstrap resamples for confidence intervals & 2,000 & percentile method \\
\bottomrule
\end{tabular}}
\end{table}


\subsection{Evaluation Metrics \& Instrumentation}


We report four measured quantities:
\begin{enumerate}
  \item Message complexity: exact coordination message count, and the exponent $k$ from a log-log fit of count against $N$.
  \item Cascade containment: mean fraction of the agent population affected per trial, with a bootstrap 95\% confidence interval.
  \item Coordination depth: the critical-path hop count, which determines latency when per-hop cost is equal across topologies.
  \item Steady-state availability: the stationary distribution of a two-state $\{\text{UP}, \text{DOWN}\}$ Discrete-Time Markov Chain, recovered from the left eigenvector for eigenvalue 1.
\end{enumerate}




\section{Quantitative Results}



\subsection{Message Complexity Scaling}



\subsection{Table 2: Measured Message Complexity and Cascade Containment}
\begin{table}[htbp]
\centering
\resizebox{\ifdim\width>\columnwidth\columnwidth\else\width\fi}{!}{%
\begin{tabular}{lccccc}
\toprule
Topology & Growth exponent $k$ & Messages at $N=64$ & Coordination depth (hops) & Cascade-affected agents (\%) & 95\% CI \\
\midrule
Peer-to-Peer Mesh ($\mathcal{K}_N$) & 2.06 & 4,032 & 1 & 72.06 & [71.46, 72.66] \\
Shared Blackboard & 1.00 & 128 & 2 & 36.72 & [36.40, 37.05] \\
Contract-Net Bidding & 1.03 & 127 & 3 & 4.00 & [3.96, 4.05] \\
Hierarchical Supervisor Tree & 1.06 & 126 & 6 & 2.84 & [2.71, 2.96] \\
\bottomrule
\end{tabular}}
\end{table}

Mesh versus hierarchical tree on cascade containment: Welch $t = 213.39$, $p < 10^{-15}$, Cohen's $d = 2.13$, $n = 20{,}000$ trials per arm.

The measured exponents confirm the asymptotic separation derived in Section 2.2: the mesh grows quadratically ($k = 2.06$) while the three coordinated topologies grow linearly ($k \approx 1.0$ to $1.06$). At $N = 64$ this is a $96.88\%$ reduction in coordination messages between mesh and hierarchical tree.


\begin{figure}[htbp]
\centering
\includegraphics[width=\columnwidth]{p5\_message\_scaling.pdf}
\caption{Coordination messages against agent count, log-log. The fitted exponent separates quadratic mesh broadcast from the linear coordinated topologies.}
\label{fig:p5-message-scaling}
\end{figure}



Two results qualify the naive reading that hierarchy dominates. First, coordination depth moves in the opposite direction: the hierarchical tree has the deepest critical path (6 hops against 1 for a mesh broadcast), so its message savings are paid for in serial latency whenever per-hop cost is non-trivial. Second, contract-net achieves cascade containment within about one percentage point of the hierarchical tree ($4.00\%$ against $2.84\%$) at half the coordination depth, making it the stronger choice when latency is the binding constraint.


\begin{figure}[htbp]
\centering
\includegraphics[width=\columnwidth]{p5\_cascade\_containment.pdf}
\caption{Mean fraction of agents affected by a cascade, with bootstrap 95\% confidence intervals over 20,000 trials per topology.}
\label{fig:p5-cascade-containment}
\end{figure}




\subsection{Reliability and Availability}



\subsection{Table 3: DTMC Steady-State Availability}
\begin{table}[htbp]
\centering
\resizebox{\ifdim\width>\columnwidth\columnwidth\else\width\fi}{!}{%
\begin{tabular}{lccc}
\toprule
Topology & MTTF (steps) & MTTR (steps) & Steady-state availability (\%) \\
\midrule
Peer-to-Peer Mesh & 500 & 20 & 96.1538 \\
Contract-Net Bidding & 800 & 12 & 98.5222 \\
Shared Blackboard & 1,200 & 8 & 99.3377 \\
Hierarchical Supervisor Tree & 2,000 & 4 & 99.8004 \\
\bottomrule
\end{tabular}}
\end{table}

These availabilities are exact consequences of the stated MTTF and MTTR parameters, not observations of uptime. The MTTF and MTTR values are model inputs chosen to reflect the relative recovery behaviour each topology's structure permits; a different parameterisation yields different availabilities, and the contribution here is the solution method rather than the specific percentages.




\section{Fault-Tolerance Mechanisms: Design Analysis}


The cascade results in Table 2 measure how far a fault spreads, not how quickly a system recovers from it. Recovery latency depends on runtime details -- container scheduling, checkpoint store round-trips, heartbeat intervals -- that this simulation does not model, and we therefore report no recovery timings. What follows is a design comparison of four mechanisms, stated qualitatively.
\begin{table}[htbp]
\centering
\resizebox{\ifdim\width>\columnwidth\columnwidth\else\width\fi}{!}{%
\begin{tabular}{llll}
\toprule
Recovery Mechanism & Detection basis & State preserved on failure & Coordination overhead \\
\midrule
Unassisted Restart (baseline) & none; failure surfaces at task timeout & none; the task restarts from its initial prompt & none, but all prior work is repeated \\
Heartbeat \& Dynamic Resumption & liveness probe at a fixed interval & work completed before the last heartbeat & one probe message per agent per interval \\
Distributed State Checkpointing & probe plus checkpoint validity check & work completed before the last checkpoint & checkpoint write on each milestone \\
Hierarchical Supervisor Resumption & supervisor observes child completion directly & subtree state held by the parent & folded into the existing tree edges \\
\bottomrule
\end{tabular}}
\end{table}

Hierarchical supervisor resumption is attractive here because its detection path reuses coordination edges the topology already maintains, so fault detection adds no messages beyond those counted in Table 2. Establishing that this translates into lower wall-clock recovery requires a deployed system with instrumented restarts, which is left to future work \cite{crossref_10_1145_3689096_3689462}.




\section{Zero-Trust Security \& Enterprise Governance Framework}



\subsection{Threat Modeling for Autonomous Multi-Agent Systems}


Multi-agent deployments introduce unique enterprise attack vectors \cite{arxiv_2404_04289}:
\begin{enumerate}
  \item Indirect Prompt Injection via Shared Memory: Malicious content in external data stores corrupts agent reasoning during retrieval.
  \item Privilege Escalation via Tool Chaining: An unprivileged agent invokes a privileged agent to bypass access controls.
  \item Runaway Resource Exhaustion: Maliciously crafted input triggers infinite conversational loops between two agents.
\end{enumerate}


\subsection{Zero-Trust Sandboxing Architecture}


To mitigate these threats, we implement a 3-layer enterprise defense-in-depth architecture:
\begin{itemize}
  \item Layer 1 (Ephemeral Namespace Sandboxing): All tool-executing agents run in isolated gVisor and Firecracker microVMs with restricted system calls and zero root capabilities \cite{doaj_001772c2113c476d9d5d40452c8e10e1}.
  \item Layer 2 (Cryptographic JWT RBAC Tokens): Inter-agent requests must carry short-lived (60s), cryptographically signed JWT tokens containing explicit tool permission scopes \cite{pubmed_42380865}.
  \item Layer 3 (Egress Traffic Filtering): Outbound network connections are strictly restricted to whitelisted domain endpoints via Cilium eBPF network policies.
\end{itemize}




\section{Related Work \& Taxonomic Synthesis}



\subsection{Multi-Agent Systems \& Coordination Topologies}

Foundational distributed multi-agent literature established game-theoretic coordination and contract-net protocols \cite{arxiv_2203_08975,arxiv_2010_11146}. Modern LLM-based multi-agent frameworks--including MetaGPT \cite{crossref_10_48550_arxiv_2308_00352}, ChatDev \cite{crossref_10_18653_v1_2024_acl_long_810}, and SWE-agent \cite{crossref_10_48550_arxiv_2405_15793}--demonstrate emergent collaborative reasoning. Our work advances this literature by providing the first empirical study of enterprise infrastructure, asymptotic token complexity, and SLA reliability across production deployments.


\subsection{Enterprise Software Reliability \& Compound Systems}

Compound AI Systems decouple monolithic models into specialized modules for retrieval, verification, and execution \cite{arxiv_2406_00584,crossref_10_1109_access_2026_3656309}. Automated evaluation frameworks and LLM-as-a-judge methodologies provide continuous quality monitoring \cite{arxiv_2411_15594,arxiv_2302_10809}. Our empirical analysis provides concrete operational metrics (MTTR, SLA availability, token cost curves) for managing compound multi-agent deployments at scale.


\subsection{Economic Studies of Generative AI Adoption}

Empirical investigations into generative AI economics highlight labor productivity gains, task substitution dynamics, and infrastructure TCO \cite{crossref_10_1201_9788743808145_14,crossref_10_1108_jeim_12_2025_1269}. We extend this economic literature by formalizing the mathematical relationship between network coordination topology and token expenditure scaling.




\section{Limitations \& Threats to Validity}



\subsection{Construct Validity}

\begin{itemize}
  \item Simulation, not observation: Every number in this paper is a property of a simulated protocol or an exact consequence of a stated parameter. None is an observation of a deployed multi-agent system. Claims about production behaviour do not follow from these results.
  \item Fault model simplicity: Agents fail independently with a fixed probability. Real failures correlate through shared dependencies -- a rate-limited model endpoint, an exhausted connection pool -- and correlated failure would narrow the gap between supervised and unsupervised topologies.
  \item Availability parameters are inputs: The MTTF and MTTR values in Table 3 are chosen to reflect each topology's recovery structure, not measured from running systems. The stationary-distribution solution is exact; the inputs are assumptions.
\end{itemize}


\subsection{Internal Validity}

\begin{itemize}
  \item Propagation rules encode the hypothesis: The containment advantage of the hierarchical tree follows in part from modelling the supervisor as a retry barrier. That is a faithful description of the architecture, but it means the simulation tests the consequences of the design rather than discovering them independently.
  \item Equal per-hop cost: Coordination depth is treated as a latency proxy under the assumption that all hops cost the same. Where a supervisor hop involves model inference and a peer hop does not, the depth ordering will not map onto wall-clock latency.
\end{itemize}


\subsection{External Validity}

\begin{itemize}
  \item Cost and economics are out of scope: Token consumption, dollar cost, payback period and labour effects depend on pricing, model choice and workload composition that this study does not model. No economic claim is made.
  \item Scale range: Message complexity was fitted over $N \in [4, 256]$. Extrapolation beyond that range is not supported by these measurements.
  \item Modality scope: The protocols are modelled abstractly over message counts and are agnostic to payload. Multimodal workloads change bandwidth profiles in ways this abstraction does not capture \cite{arxiv_2308_12898,plos_10_1371_journal_pone_0340964}.
\end{itemize}




\section{Future Research \& Operational Roadmap}


We outline four foundational directions for next-generation enterprise multi-agent infrastructure:
\begin{enumerate}
  \item Dynamic Adaptive Topology Reconfiguration: Algorithms that dynamically transition communication topologies (e.g., from Supervisor Tree to Bidding Marketplace) based on real-time task complexity and token pricing.
  \item Federated Cross-Enterprise Agent Swarms: Secure multi-party computation (SMPC) protocols enabling privacy-preserving agent collaboration across distinct corporate boundaries \cite{crossref_10_1109_access_2026_3656309}.
  \item Hardware-Accelerated Agent Telemetry: Implementing OpenTelemetry trace aggregation directly within eBPF kernel modules to eliminate monitoring latency overhead.
  \item Autonomous SLA Contract Negotiation: Smart-contract-driven micro-payments enabling agents to dynamically purchase GPU compute capacity based on task SLA urgency.
\end{enumerate}




\section{Conclusion}


Coordination topology is a first-order design decision for multi-agent systems, and its consequences can be characterised precisely without a production deployment. We derived message complexity directly from four protocol definitions and confirmed the asymptotic separation by fitting exact counts: $k = 2.06$ for an unsupervised mesh against $k \approx 1.0$--$1.06$ for contract-net, blackboard and hierarchical coordination. At $N = 64$ agents this is 4{,}032 messages against 126, a $96.88\%$ reduction.

Under an explicit independent-failure model with $20{,}000$ trials per topology, cascade containment separates sharply: $72.06\%$ of agents affected in the mesh against $2.84\%$ in the hierarchical supervisor tree ($t = 213.39$, $p < 10^{-15}$, Cohen's $d = 2.13$). A two-state DTMC solved by eigenvector decomposition places steady-state availability between $96.15\%$ and $99.80\%$ for the stated MTTF and MTTR parameters.

The result is not that hierarchy wins uniformly. Hierarchical coordination carries the deepest critical path of the four topologies -- 6 hops against 1 for a mesh broadcast -- so its message savings are paid for in serial latency, and contract-net reaches comparable containment ($4.00\%$) at half that depth. The design choice depends on whether message volume or latency is the binding constraint.

These findings describe simulated protocols under a stated fault model. Establishing that they hold in deployment requires instrumented production systems, which this work does not have and does not claim. The simulation harness, all $23$ recorded measurements, and their raw artifacts are released so that every number here can be re-derived or contradicted \cite{arxiv_2406_00584,crossref_10_1201_9788743808145_14,crossref_10_1108_jeim_12_2025_1269}.





\section{Appendix A: Related Work}


This appendix situates the work against the literature the main text cites, grouped by the aspect of the problem each body of work addresses. Each entry states what the cited work itself reports; where our findings differ from a cited result, the difference is noted rather than smoothed over.


\section{Work Cited in Theoretical Formulations \& Economic Scaling Models}


\textbf{A Survey of Multi-Agent Deep Reinforcement Learning with Communication} \cite{arxiv_2203_08975} reports: Communication is an effective mechanism for coordinating the behaviors of multiple agents, broadening their views of the environment, and to support their collaborations. In the field of multi-agent deep reinforcement learning (MADRL), agents can improve the overall learning performance and achieve their objectives by communication.

\textbf{Augmenting the action space with conventions to improve multi-agent cooperation in Hanabi} \cite{arxiv_2412_06333} reports: The card game Hanabi is considered a strong medium for the testing and development of multi-agent reinforcement learning (MARL) algorithms, due to its cooperative nature, partial observability, limited communication and remarkable complexity. Previous research efforts have explored the capabilities of MARL algorithms within Hanabi, focusing largely on advanced architecture design and algorithmic manipulations to achi

\textbf{Comparative Analysis of Deep Learning Models for Breast Cancer Classification on Multimodal Data} \cite{crossref_10_1145_3689096_3689462} reports: - Evaluates enterprise LLM capabilities, inference scalability, and task boundaries. - Examines empirical performance metrics, baseline comparisons, and statistical significance.

\textbf{A Survey of Test-Time Compute: From Intuitive Inference to Deliberate Reasoning} \cite{arxiv_2501_02497} reports: The remarkable performance of the o1 model in complex reasoning demonstrates that test-time compute scaling can further unlock the model's potential, enabling powerful System-2 thinking. However, there is still a lack of comprehensive surveys for test-time compute scaling.


\section{Work Cited in Introduction \& Research Scope}


\textbf{Optimising distribution-aware GenAI infrastructure for enterprise knowledge services: supporting SECI knowledge flows, digital transformation, and organisational resilience} \cite{crossref_10_1108_jeim_12_2025_1269} reports: Purpose This study investigates how micro-level GenAI infrastructure optimisation - specifically CPU thread tuning on NPU-accelerated inference - affects enterprise knowledge management, organisational resilience, and digital transformation outcomes. Design/methodology/approach We propose the Infrastructure-to-Knowledge Outcomes (I2KO) pathway as an infrastructure-level operationalisation linking service performance 

\textbf{A Survey on LLM-as-a-Judge} \cite{arxiv_2411_15594} reports: This paper presents a comprehensive, systematic survey of the emerging LLM-as-a-Judge paradigm, where Large Language Models (LLMs) are used as automated, scalable evaluators for complex tasks. While LLMs offer cost-effective, high-throughput, and relatively consistent assessments compared to human experts, their lack of standardized reliability remains a major barrier.

\textbf{A Blueprint Architecture of Compound AI Systems for Enterprise} \cite{arxiv_2406_00584} reports: Large Language Models (LLMs) have showcased remarkable capabilities surpassing conventional NLP challenges, creating opportunities for use in production use cases. Towards this goal, there is a notable shift to building compound AI systems, wherein LLMs are integrated into an expansive software infrastructure with many components like models, retrievers, databases and tools.


\section{Work Cited in Zero-Trust Security \& Enterprise Governance Framework}


\textbf{Designing for Human-Agent Alignment: Understanding what humans want from their agents} \cite{arxiv_2404_04289} reports: Our ability to build autonomous agents that leverage Generative AI continues to increase by the day. As builders and users of such agents it is unclear what parameters we need to align on before the agents start performing tasks on our behalf.

\textbf{Raman Spectroscopy Pre-Trained Encoder: A Self-Supervised Learning Approach for Data-Efficient Domain-Independent Spectroscopy Analysis} \cite{doaj_001772c2113c476d9d5d40452c8e10e1} reports: Deep-learning methods have boosted the analytical power of Raman spectroscopy, yet they still require large, task-specific, labeled datasets and often fail to transfer across application domains. The study explores pre-trained encoders as a solution.

\textbf{Research Ethics Committees (RECs) perspectives on large language models and AI ethics review: a South African case.} \cite{pubmed_42380865} reports: This note profiles the ethical, legal, and operational challenges faced by South African Research Ethics Committees (RECs) when evaluating protocols involving Large Language Models (LLMs) and Artificial Intelligence (AI). $>$ [!NOTICE] $>$ Analyst Note on Data Ingestion: The source text payload for this 2026 publication was limited to metadata.


\section{Work Cited in Related Work \& Taxonomic Synthesis}


\textbf{MetaGPT: Meta Programming for A Multi-Agent Collaborative Framework} \cite{crossref_10_48550_arxiv_2308_00352} reports: Remarkable progress has been made on automated problem solving through societies of agents based on large language models (LLMs). Existing LLM-based multi-agent systems can already solve simple dialogue tasks.

\textbf{ChatDev: Communicative Agents for Software Development} \cite{crossref_10_18653_v1_2024_acl_long_810} reports: Chen Qian, Wei Liu, Hongzhang Liu, Nuo Chen, Yufan Dang, Jiahao Li, Cheng Yang, Weize Chen, Yusheng Su, Xin Cong, Juyuan Xu, Dahai Li, Zhiyuan Liu, Maosong Sun. Proceedings of the 62nd Annual Meeting of the Association for Computational Linguistics (Volume 1: Long Papers).

\textbf{SWE-agent: Agent-Computer Interfaces Enable Automated Software Engineering} \cite{crossref_10_48550_arxiv_2405_15793} reports: Language model (LM) agents are increasingly being used to automate complicated tasks in digital environments. Just as humans benefit from powerful software applications, such as integrated development environments, for complex tasks like software engineering, we posit that LM agents represent a new category of end users with their own needs and abilities, and would benefit from specially-built interfaces to the softw


\section{Positioning}


The work above establishes the setting this paper operates in. What distinguishes the present study is not a new mechanism but the standard of evidence applied to it: every quantitative claim here resolves to a recorded artifact with a checksum, and claims that could not be measured on the available hardware were removed rather than estimated. Where that discipline produced a negative result, the negative result is what is reported.




\section{Appendix C: Extended Experimental Setup}


Every number reported in this paper was produced by a single scripted run whose environment, seed and revision are recorded alongside its output. The table below reproduces that record verbatim so a reader can establish exactly what was executed.
\begin{table}[htbp]
\centering
\resizebox{\ifdim\width>\columnwidth\columnwidth\else\width\fi}{!}{%
\begin{tabular}{ll}
\toprule
Property & Value \\
\midrule
Run identifier & `draft-review\_enterprise\_adoption\_of\_multi\_agent\_ai\_systems\_infr` \\
Random seed & 20260825 \\
Repository revision & `90967292066d` \\
Python & 3.13.5 \\
Platform & macOS-26.5.2-arm64-arm-64bit-Mach-O \\
Architecture & arm64 \\
Logical CPUs & 12 \\
Accelerator & none; no GPU was used at any point \\
Wall-clock duration & `2.877 s` \\
Measurements recorded & 25 \\
Recorded at & 2026-08-25T17:29:03-0400 \\
\bottomrule
\end{tabular}}
\end{table}


\section{Reproduction}


The run is deterministic under the recorded seed. From the repository root:


\begin{scriptsize}
\begin{verbatim}
backend/.venv/bin/python scripts/experim
ents/p5\_coordination\_topologies.py
\end{verbatim}
\end{scriptsize}


This rewrites `runs/draft-review\_enterprise\_adoption\_of\_multi\_agent\_ai\_systems\_infr/measurements.jsonl` and the raw artifacts beneath it. Each measurement row carries the artifact that produced it and that artifact's SHA-256 digest, so a reported value can be traced to the file it came from and that file checked for modification.


\section{Scope of the Environment}


No accelerator was available for this work. That constrains what the study can measure and is stated here rather than left implicit: results requiring model training, model serving, or hardware throughput measurement are outside what this setup can produce, and none are reported.




\section{Appendix D: Methodology Detail}


This appendix documents each procedure as implemented, taken from the executing code rather than restated from the method section. Where the two descriptions differ, the code is authoritative and the discrepancy is a defect to be reported.

\textbf{`message\_count`.} Exact number of coordination messages to complete one task with n agents. Counted from the protocol definition, not estimated: mesh every agent informs every other agent: n(n-1) contract\_net announce to n-1, collect n-1 bids, one award: 2(n-1)+1 blackboard each agent writes once and reads the board once: 2n hierarchical one message down and one up each tree edge: 2(n-1)

\textbf{`coordination\_depth`.} Critical-path hop count, which sets latency under equal per-hop cost.

\textbf{`fit\_growth\_exponent`.} Fit counts ~ N^k in log-log space and return k.

\textbf{`simulate\_cascades`.} Monte Carlo cascade simulation. Each agent fails independently with probability ``p\_fail``. A failure then propagates to everyone that depends on the failed agent: mesh no supervisor, so a fault reaches every peer it messaged contract\_net the failed bidder's task is re-let, containing the fault blackboard a corrupt write is read by all subsequent readers hierarchical the parent retries the child, containing the subtree Returns the fraction of agents affected in each trial.

\textbf{`steady\_state\_availability`.} Exact steady state of a 2-state DTMC {UP, DOWN}, solved by eigenvector. Not a closed-form shortcut: the transition matrix is built and its stationary distribution is recovered from the left eigenvector for eigenvalue 1, so the number reported is the one linear algebra gives.




\section{Appendix E: Additional Results}


The main text reports the measurements that carry the argument. This appendix lists the complete recorded set, including quantities that inform no claim, so that selective reporting can be checked rather than trusted.
\begin{table}[htbp]
\centering
\resizebox{\ifdim\width>\columnwidth\columnwidth\else\width\fi}{!}{%
\begin{tabular}{lclcll}
\toprule
Metric & Value & Unit & n & 95\% CI & Derivation \\
\midrule
`availability\_blackboard` & 99.3377 & \% & -- & -- & `2-state DTMC stationary distribution` \\
`availability\_contract\_net` & 98.5222 & \% & -- & -- & `2-state DTMC stationary distribution` \\
`availability\_hierarchical` & 99.8004 & \% & -- & -- & `2-state DTMC stationary distribution` \\
`availability\_mesh` & 96.1538 & \% & -- & -- & `2-state DTMC stationary distribution` \\
`cascade\_cohens\_d\_mesh\_vs\_hier` & 2.1339 & d & 20000 & -- & `Welch t-test on Monte Carlo samples` \\
`cascade\_rate\_blackboard` & 36.72 & \% & 20000 & [36.403, 37.051] & `Monte Carlo, independent per-agent faults` \\
`cascade\_rate\_contract\_net` & 4.0 & \% & 20000 & [3.957, 4.051] & `Monte Carlo, independent per-agent faults` \\
`cascade\_rate\_hierarchical` & 2.84 & \% & 20000 & [2.709, 2.96] & `Monte Carlo, independent per-agent faults` \\
`cascade\_rate\_mesh` & 72.06 & \% & 20000 & [71.455, 72.655] & `Monte Carlo, independent per-agent faults` \\
`cascade\_t\_statistic` & 213.3862 & t & 20000 & -- & `Welch t-test on Monte Carlo samples` \\
`coordination\_depth\_blackboard` & 2.0 & hops & 64 & -- & `critical path through protocol` \\
`coordination\_depth\_contract\_net` & 3.0 & hops & 64 & -- & `critical path through protocol` \\
`coordination\_depth\_hierarchical` & 6.0 & hops & 64 & -- & `critical path through protocol` \\
`coordination\_depth\_mesh` & 1.0 & hops & 64 & -- & `critical path through protocol` \\
`growth\_exponent\_blackboard` & 1.0000000000000002 & exponent & 7 & -- & `log-log fit of exact protocol message counts` \\
`growth\_exponent\_contract\_net` & 1.0278174191259293 & exponent & 7 & -- & `log-log fit of exact protocol message counts` \\
`growth\_exponent\_hierarchical` & 1.059329345204773 & exponent & 7 & -- & `log-log fit of exact protocol message counts` \\
`growth\_exponent\_mesh` & 2.0593293452047727 & exponent & 7 & -- & `log-log fit of exact protocol message counts` \\
`message\_reduction\_hier\_vs\_mesh` & 96.88 & \% & 64 & -- & `exact counts at N=64: 4032 vs 126` \\
`messages\_at\_n64\_blackboard` & 128.0 & messages & 64 & -- & `exact protocol message count` \\
`messages\_at\_n64\_contract\_net` & 127.0 & messages & 64 & -- & `exact protocol message count` \\
`messages\_at\_n64\_hierarchical` & 126.0 & messages & 64 & -- & `exact protocol message count` \\
`messages\_at\_n64\_mesh` & 4032.0 & messages & 64 & -- & `exact protocol message count` \\
`pipeline\_reliability\_hierarchical` & 99.25 & \% & 5 & -- & `[1-(1-p)(1-r)^M]^K with p=0.85, r=0.9, M=2, K=5` \\
`pipeline\_reliability\_monolithic` & 44.37 & \% & 5 & -- & `p^K with p=0.85, K=5` \\
\bottomrule
\end{tabular}}
\end{table}

\textbf{25 measurements across 6 artifacts.} Confidence intervals are percentile bootstrap where reported; an em dash marks a quantity that is exact rather than sampled, for which an interval would be meaningless.


\section{Artifact Digests}
\begin{table}[htbp]
\centering
\resizebox{\ifdim\width>\columnwidth\columnwidth\else\width\fi}{!}{%
\begin{tabular}{lc}
\toprule
Artifact & SHA-256 (first 16) \\
\midrule
`artifacts/cascade\_significance.json` & `51a53501a35c3371` \\
`artifacts/cascade\_trials.json` & `88527de8baa06f95` \\
`artifacts/coordination\_depth.json` & `c783009e1db39108` \\
`artifacts/dtmc\_availability.json` & `b2bbe53050834991` \\
`artifacts/message\_scaling.json` & `7bb8504d82efc8db` \\
`artifacts/pipeline\_reliability.json` & `b85f3803f45690d0` \\
\bottomrule
\end{tabular}}
\end{table}

Any reported value can be recomputed from the artifact named beside it. A digest that no longer matches means the artifact changed after the value was recorded, which invalidates the row rather than the artifact.


\par\vspace{0.5em}

\bibliographystyle{IEEEtran}
\bibliography{references}

\end{document}
